% =====================================================================
%  MARCO REFERENCIAL
%
%  Los §2.1 y §2.2 se trasvasan del informe de avance del marco teórico
%  (Documentos/Entregables/Actualizacion_Marco_Teorico_Conceptual.tex),
%  adaptados a la jerarquía de capítulo/sección/subsección del documento final.
%
%  Los §2.3 y §2.4 son contenido nuevo: no existían en ningún documento del
%  proyecto. Llevan marcados en comentarios los puntos que hay que VERIFICAR
%  antes de entregar; no se afirma nada normativo que no esté comprobado.
% =====================================================================
\chapter{Marco referencial}

Este capítulo reúne el vocabulario técnico, los fundamentos teóricos y el
encuadre normativo y ambiental con que opera el proyecto. Se organiza en cuatro
marcos. El conceptual fija el significado de los términos que el resto del
documento usa sin volver a definirlos. El teórico expone las decisiones de
arquitectura y sus fundamentos, incluidas aquellas que la implementación obligó
a revisar. El legal sitúa el trabajo frente a la normativa aplicable, y el
ambiental frente a las consideraciones de consumo y ciclo de vida del hardware
empleado.

% =====================================================================
\section{Marco conceptual}
% =====================================================================

Se definen a continuación los términos técnicos necesarios para la lectura del
documento, considerando un público con formación en ingeniería electrónica.
La organización en seis categorías se conserva del anteproyecto; dentro de cada
una, los términos que la implementación introdujo se añaden a los ya definidos.

\subsection{Navegación y percepción}

\begin{description}[style=nextline,leftmargin=1.5em]
  \item[Navegación autónoma] Desplazamiento sin intervención humana que integra
    percepción, localización y control.
  \item[Odometría] Estimación de posición por integración de mediciones de
    actuadores o de sensores exteroceptivos.
  \item[Zona de transición] Punto crítico para el cambio de nivel en
    edificaciones con múltiples pisos.
  \item[Mapa de costos (\textit{costmap})] Representación del entorno usada para
    la planificación, que asigna a cada celda un costo de tránsito en lugar de un
    simple estado libre u ocupado. Existe una capa global, que cubre el mapa
    completo, y una local, limitada al entorno inmediato del robot.
  \item[AMCL (\textit{Adaptive Monte Carlo Localization})] Algoritmo de
    localización basado en un filtro de partículas que estima la posición y
    orientación del robot combinando la odometría con las lecturas del sensor
    láser contra un mapa conocido.
  \item[Verificador de meta (\textit{goal checker})] Componente que decide si el
    robot alcanzó su objetivo, comparando posición y, opcionalmente, orientación
    contra una tolerancia.
  \item[Observabilidad de un estado] Propiedad de una variable física de poder
    estimarse a partir de las mediciones disponibles. Una variable puede ser
    matemáticamente inobservable en una configuración concreta de entorno y
    sensor, sin que exista error de calibración que lo corrija.
  \item[Modelo de movimiento Reeds-Shepp] Conjunto de trayectorias óptimas para
    un vehículo con curvatura acotada que puede moverse hacia adelante y hacia
    atrás, usado por el planificador híbrido para generar maniobras de inversión
    de marcha.
  \item[Cúspide] Punto de una trayectoria planificada en el que el vehículo
    invierte el sentido de marcha.
\end{description}

\subsection{Control y locomoción}

\begin{description}[style=nextline,leftmargin=1.5em]
  \item[Cinemática de robots móviles] Modelo matemático del movimiento en
    función de las velocidades de los actuadores.
  \item[Modelo cinemático tipo automóvil (dirección Ackermann)] Representación
    geométrica y matemática de un vehículo de cuatro ruedas con dirección en el
    eje delantero, como la plataforma empleada. Relaciona el ángulo de giro del
    servomotor con el radio de giro del vehículo, asegurando que las ruedas
    interior y exterior giren en torno a un mismo centro instantáneo de rotación.
  \item[PWM (\textit{Pulse Width Modulation})] Técnica para controlar motores o
    servomotores mediante el ciclo de trabajo de una señal.
  \item[Control PID] Algoritmo de control por retroalimentación
    proporcional-integral-derivativo.
  \item[Planificador híbrido A* (\textit{Smac Hybrid-A*})] Algoritmo de búsqueda
    de rutas que opera sobre el espacio continuo de posición y orientación,
    respetando las restricciones cinemáticas del vehículo, a diferencia de un
    planificador de cuadrícula que ignora la orientación.
  \item[Controlador de seguimiento regulado (\textit{Regulated Pure Pursuit})]
    Algoritmo de control que persigue un punto móvil sobre la trayectoria
    planificada, regulando su velocidad según la curvatura y la proximidad a
    obstáculos.
  \item[Banda muerta de tracción] Intervalo de velocidades comandadas por debajo
    del cual el vehículo no llega a moverse, porque la señal resultante no vence
    el rozamiento estático.
\end{description}

\subsection{Comunicación y sistemas distribuidos}

\begin{description}[style=nextline,leftmargin=1.5em]
  \item[Middleware robótico] Capa de software para la comunicación entre módulos
    del sistema.
  \item[ROS (\textit{Robot Operating System})] Middleware con nodos, tópicos de
    publicación y suscripción, servicios de petición y respuesta, y acciones
    \cite{macenski2022ros2}.
  \item[DDS (\textit{Data Distribution Service})] Protocolo de transporte de
    publicación y suscripción sobre el que ROS 2 implementa sus tópicos,
    servicios y acciones, con descubrimiento automático de participantes.
  \item[Dominio ROS] Identificador numérico que segmenta el descubrimiento DDS.
    Dos procesos con distinto dominio no se ven entre sí salvo configuración
    explícita.
  \item[Espacio de nombres (\textit{namespace})] Prefijo aplicado a los tópicos,
    servicios y acciones de un nodo, usado en este trabajo para aislar a cada
    agente dentro de un mismo dominio DDS.
  \item[Acción ROS] Patrón de comunicación asíncrona para tareas de larga
    duración, compuesto por un objetivo, retroalimentación periódica y un
    resultado final, con posibilidad explícita de cancelación.
  \item[Servidor y cliente de acción] Nodo que ofrece una acción y nodo que la
    invoca. En este sistema el coordinador es cliente de la acción de navegación
    ante cada agente y servidor de la acción de guiado ante la interfaz.
  \item[Callback de cancelación] Función que un servidor de acción debe registrar
    explícitamente para poder aceptar solicitudes de cancelación; sin ella, el
    comportamiento por omisión es rechazarlas.
  \item[Calidad de servicio (QoS)] Conjunto de políticas de entrega de un tópico
    o servicio, configurable de forma independiente para cada canal.
  \item[Puente \textit{rosbridge}] Servidor que traduce entre el protocolo nativo
    de ROS 2 y un protocolo JSON sobre WebSocket, permitiendo que un cliente web
    sin dependencias de ROS publique, se suscriba y opere acciones.
\end{description}

\subsection{Arquitectura del sistema}

\begin{description}[style=nextline,leftmargin=1.5em]
  \item[Asignación de tareas] Distribución de actividades entre robots según su
    disponibilidad y ubicación \cite{gerkey2004taxonomy}.
  \item[Protocolo de relevo] Transferencia secuencial de responsabilidad entre
    agentes en puntos definidos.
  \item[Máquina de estados de misión] Conjunto de valores discretos y ordenados
    que describen el progreso de una misión, publicados para que tanto la
    interfaz como el registro de evaluación conozcan el punto de ejecución.
  \item[Planificador puro] Componente que calcula la secuencia de tramos de una
    misión a partir del catálogo de puntos y del par origen-destino, sin depender
    del middleware ni de ningún estado externo, lo que permite verificarlo con
    cientos de casos en milisegundos y sin simulador.
  \item[Punto de transferencia] Punto de interés marcado en el catálogo como el
    lugar físico, típicamente una escalera, en el que un agente entrega y otro
    recibe la continuidad del guiado al cambiar de nivel.
  \item[Registro de misión] Documento estructurado, compuesto a partir de la
    grabación de la misión y validado contra un esquema formal, que constituye la
    evidencia primaria de una ejecución para efectos de la evaluación
    experimental.
\end{description}

\subsection{Interacción humano-robot}

\begin{description}[style=nextline,leftmargin=1.5em]
  \item[Interfaz gráfica de usuario (GUI)] Entorno visual para la interacción
    humano-máquina.
  \item[HRI (\textit{Human-Robot Interaction})] Diseño y evaluación de sistemas
    robóticos destinados a interactuar con personas \cite{jimenez2024arquitectura}.
  \item[Solicitud de servicio] Petición de un usuario para recibir asistencia de
    guiado.
  \item[Interfaz responsiva de selección] Componente que permite ubicar un punto
    de interés por nombre mediante búsqueda incremental, sin perder la capacidad
    de mostrar el catálogo completo.
  \item[Retroalimentación literal] Principio de diseño según el cual el texto
    mostrado al usuario se redacta íntegramente en el componente que conoce el
    estado del sistema y se presenta sin reinterpretación en la interfaz.
\end{description}

\subsection{Métricas de evaluación}

\begin{description}[style=nextline,leftmargin=1.5em]
  \item[Tiempo de respuesta] Intervalo entre la solicitud del usuario y el inicio
    de la tarea por el robot.
  \item[Tiempo de asignación] Intervalo entre la aceptación de la solicitud y la
    designación del agente que la atenderá.
  \item[Tasa de éxito] Proporción de solicitudes completadas exitosamente.
  \item[Continuidad de servicio] Capacidad de mantener la asistencia durante las
    transiciones entre niveles.
  \item[Intervalo de confianza de Wilson] Método para estimar el intervalo de
    confianza de una proporción, más confiable que la aproximación normal cuando
    el tamaño de muestra es pequeño o la proporción se acerca a 0 o a 1.
  \item[Falsabilidad de un criterio experimental] Propiedad de un criterio de
    aceptación de poder, en principio, arrojar un resultado negativo bajo alguna
    ejecución posible del sistema. Un criterio que no puede fallar no aporta
    evidencia, sin importar cuántas veces se verifique satisfactoriamente.
  \item[Invariante estructural] Propiedad que se cumple siempre por el propio
    diseño del sistema, y que por tanto no debe presentarse como un resultado
    experimental medido, sino declararse como tal.
  \item[Cota superior por misión] Forma de reportar una métrica cuando el
    instrumento de medición no tiene resolución suficiente para dar un valor
    exacto, reportando el límite superior que sí puede garantizar.
\end{description}

% =====================================================================
\section{Marco teórico}
% =====================================================================

\subsection{Robot móvil, sistema multi-robot y navegación}

El sistema es un conjunto de dos robots móviles terrestres que operan como
agentes colaborativos, cada uno dedicado a un nivel del edificio, y que en
conjunto conforman un sistema multi-robot \cite{varghese2023multirobot}. La tarea
que ejecutan, guiar a una persona desde un punto de origen hasta un destino con
selección de la tarea a través de una interfaz móvil, constituye interacción
humano-robot \cite{jimenez2024arquitectura}.

La navegación de cada agente se apoya en tres pilares: la localización, que
determina la posición y orientación del robot respecto al entorno; la
planificación de trayectorias, que calcula una ruta factible respetando las
restricciones físicas del vehículo; y el control de movimiento, que genera las
señales necesarias para ejecutar esa ruta. Los tres se alimentan de percepción
continua del entorno, provista por un sensor láser de barrido plano.

\subsection{ROS 2 como capa de integración, en dos distribuciones}

El sistema usa ROS 2 \cite{macenski2022ros2} como middleware de comunicación
entre todos sus componentes, en dos distribuciones simultáneas: una en la
simulación, porque el paquete de integración físico-simulada no está liberado
sobre distribuciones posteriores, y otra en los vehículos físicos, determinada
por la versión del sistema operativo que soportan sus tarjetas.

Esto fija una restricción que gobierna buena parte de la arquitectura: el tipo de
mensaje que constituye la única vía de mando del contrato de interfaces tiene
definición distinta entre ambas distribuciones, de modo que un coordinador
compilado en una no puede mandar directamente a un agente compilado en la otra.
\textbf{La compatibilidad de código fuente entre distribuciones de ROS 2 no
implica compatibilidad de mensajes en tiempo de ejecución}, y esa distinción no
aparece enunciada en la documentación de alto nivel del middleware. La mitigación
adoptada mantiene una variante de configuración por distribución, derivada de una
base común.

Sobre esa capa, la arquitectura de coordinación es centralizada: un único nodo
coordinador actúa como cliente de acción hacia cada agente y como servidor de
acción hacia la interfaz de usuario. En la taxonomía de asignación de tareas de
Gerkey y Matarić \cite{gerkey2004taxonomy}, corresponde a una asignación
centralizada e instantánea: el coordinador determina el agente responsable en el
momento de aceptar la solicitud, a partir del nivel de los puntos de origen y
destino, sin negociación entre agentes. Con dos agentes y una asignación
determinada por el nivel del punto solicitado, un mecanismo de subasta añadiría
complejidad sin aportar flexibilidad.

\subsection{Navegación autónoma sobre cinemática tipo automóvil}

La navegación de cada agente se ejecuta sobre el marco de navegación estándar de
ROS 2 \cite{sojan2025autonomous}. La plataforma sigue un modelo cinemático tipo
automóvil: no holonómico, incapaz de girar sobre su propio eje. La configuración
por omisión de ese marco asume un robot de tracción diferencial, de modo que dos
de sus tres capas se sustituyen:

\begin{itemize}[leftmargin=*]
  \item \textbf{Planificación global:} un planificador híbrido A* con modelo de
    movimiento Reeds-Shepp, que admite marcha atrás y genera explícitamente las
    maniobras de inversión de marcha que un vehículo no holonómico necesita para
    alcanzar orientaciones que su geometría no permite alcanzar de frente.
  \item \textbf{Control de seguimiento:} un controlador de persecución pura
    regulada con marcha atrás habilitada, que ejecuta físicamente las cúspides
    que el planificador propuso.
\end{itemize}

Esta configuración tiene un límite caracterizado con evidencia propia durante el
desarrollo. El controlador reduce dinámicamente su distancia de anticipación
cuando el vehículo se aproxima a una cúspide muy cercana al punto objetivo, y por
debajo de un umbral geométrico concreto una guarda de seguridad del propio
controlador anula la curvatura calculada: el vehículo sigue invirtiendo el
sentido de marcha, pero la dirección deja de responder, lo que produce episodios
de decenas de maniobras en el mismo punto. No es un defecto del planificador ni
de la coordinación, sino un límite del controlador de referencia operando en el
borde de su propio dominio de validez.

\subsection{Localización y mapeo: observabilidad en un corredor uniforme}

La localización se apoya en un filtro de partículas sobre el mapa del entorno, y
el mapa se construye a partir de odometría láser obtenida por correspondencia de
barridos sucesivos. En un corredor recto y de sección uniforme, la posición
\textbf{longitudinal} del robot no es observable a partir de un sensor láser
plano combinado con ese método: cada barrido sucesivo es indistinguible del
anterior a lo largo del eje del pasillo, de modo que el estimador devuelve un
desplazamiento cercano a cero y acierta desde el punto de vista de la información
disponible. La ausencia de estructura lateral variable elimina la información
necesaria para estimar el avance; el algoritmo no la pierde por error.

Esta propiedad distingue un \textbf{error de calibración}, corregible ajustando
parámetros, de una \textbf{inobservabilidad estructural}, que ningún ajuste
corrige porque la información no está en la señal. La condición de aceptación de
un mapa se define en consecuencia no como un porcentaje fijo de cobertura, sino
como la exigencia de que el entorno aporte estructura no uniforme dentro del
alcance del sensor a lo largo de todo el recorrido.

\subsection{Simulación multi-robot: aislamiento por espacio de nombres}

La simulación multi-robot se implementa con un proceso de simulación
independiente por robot y un único dominio DDS compartido entre ambos; el
aislamiento entre agentes se logra por espacio de nombres, incluido el marco de
referencia del mapa de cada robot, y no por segmentación de dominio. Un
identificador de dominio distinto por robot, que es la alternativa más directa
para aislar tópicos, bloquea el descubrimiento entre el coordinador y los
agentes salvo configuración adicional, de modo que la separación por dominios y
la comunicación entre agentes resultan mutuamente excluyentes bajo la
configuración por omisión.

\subsection{La interfaz humano-robot como frontera de protocolo}

La interfaz web se comunica con el middleware a través de un puente que expone
las acciones de ROS 2 como un protocolo JSON sobre WebSocket. Este protocolo es
distinto del que implementa la librería de cliente más difundida en el ecosistema
web de ROS, cuyo cliente de acciones corresponde a la generación anterior del
middleware y es incompatible con las acciones de ROS 2. La interfaz de este
trabajo es, por esa razón, un cliente propio escrito directamente sobre el
protocolo del puente, y sin dependencias de red externas, ya que el requisito
exige que funcione sin salida a internet.

Un segundo principio rige la interfaz: no redacta lógica de negocio. El texto que
se muestra al usuario en cada estado de misión lo redacta el coordinador y la
interfaz lo presenta literal. Es una aplicación directa del principio de
separación de responsabilidades entre lógica de dominio y capa de presentación:
si el navegador interpretara el estado numérico para decidir qué texto mostrar,
duplicaría en el cliente una lógica que ya existe en el planificador, con riesgo
de que las dos copias diverjan.

\subsection{Evaluación experimental: métricas, significancia y falsabilidad}

La evaluación se rige por cuatro métricas con definición operativa (tiempo de
respuesta, tiempo de asignación, tasa de éxito y continuidad de servicio entre
niveles), medidas sobre una campaña de misiones sorteadas, con el protocolo
fijado \textit{antes} de instrumentar el sistema. Dos principios rigen su
tratamiento.

El primero es estadístico: con tamaños de muestra del orden de treinta
observaciones, el intervalo de confianza apropiado para una proporción binomial
es el de Wilson, y no la aproximación normal, que subestima el error en ese
régimen. La diferencia no es cosmética: determina si la afirmación defendible es
una cota inferior sobre la tasa de éxito o un valor puntual, y solo la primera lo
es con esa muestra.

El segundo es metodológico: \textbf{un criterio de aceptación experimental solo
constituye evidencia si puede fallar}. Durante el desarrollo se detectó que el
criterio asociado a la variable de respuesta principal estaba definido de forma
que ninguna ejecución posible del sistema podía incumplirlo, porque los dos
únicos productores de un valor de incumplimiento ocurren, por construcción del
código, fuera de la ventana temporal en que el criterio se evalúa. Un conjunto de
verificaciones satisfactorias de un criterio que no puede resultar negativo no es
evidencia de que el sistema cumpla la propiedad; es evidencia de que el criterio
está mal planteado. La corrección adoptada no inventa una ruta de fallo
artificial, lo que equivaldría a instrumentar el experimento para que parezca más
informativo de lo que es, sino que reformula la afirmación como un invariante
estructural del diseño, respaldado por una magnitud que sí varía y sí puede
reportarse.

% =====================================================================
\section{Marco legal}
% =====================================================================

% VERIFICAR ANTES DE ENTREGAR. Las dos políticas CONPES están citadas en el
% anteproyecto y son las que el proyecto ya invocaba. Lo demás de esta sección
% se plantea como encuadre y NO afirma obligaciones concretas: para citar
% legislación de accesibilidad o de datos personales hay que comprobar la norma
% vigente y su aplicabilidad, y eso no se puede hacer de memoria.

El proyecto se desarrolla en el marco de la política pública colombiana de
transformación digital y de ciencia, tecnología e innovación. El Documento CONPES
3975 \cite{conpes3975}, que fija la Política Nacional para la Transformación
Digital e Inteligencia Artificial, y el Documento CONPES 4069 \cite{conpes4069},
que establece la Política Nacional de Ciencia, Tecnología e Innovación, instan a
la academia a desarrollar prototipos tecnológicos que apropien herramientas de la
Cuarta Revolución Industrial, entre ellas la robótica autónoma, para generar
soluciones de diseño local con impacto social. El presente trabajo se inscribe
en esa línea: un prototipo de asistencia robótica construido sobre plataformas de
bajo costo y herramientas de software abierto, dentro de las limitantes
presupuestales de un entorno universitario.

En cuanto a la presentación del propio documento, se sigue la norma NTC 1486 del
ICONTEC \cite{ntc1486}, que regula la presentación de tesis, trabajos de grado y
otros trabajos de investigación, y que la plantilla institucional declara como
su base.

% PENDIENTE DE VERIFICAR Y REDACTAR, con la norma en la mano:
%  (a) Accesibilidad en edificaciones y derechos de personas con discapacidad.
%      El proyecto se justifica socialmente por la movilidad inclusiva, así que
%      corresponde citar el marco de accesibilidad vigente. NO se redacta aquí
%      sin comprobar la norma aplicable y su versión en vigor.
%  (b) Tratamiento de datos personales. Hoy la interfaz NO registra identidad
%      del usuario -solo origen, destino y tiempos-, de modo que puede no haber
%      tratamiento de datos personales. Conviene decirlo explícitamente y
%      justificarlo, en vez de omitir el punto.
%  (c) Uso del espectro para la red inalámbrica del sistema, si aplica.

\subsection{Alcance de la operación y responsabilidad}

El alcance del proyecto excluye explícitamente la interacción mecánica con la
infraestructura del edificio: los agentes no accionan puertas ni botones de
ascensor, no transportan carga y no atraviesan las zonas de transición vertical.
Esa delimitación, adoptada por razones técnicas y de costo, tiene además una
consecuencia normativa favorable, pues evita intervenir sistemas sujetos a
regulación específica de seguridad en edificaciones.

% =====================================================================
\section{Marco ambiental}
% =====================================================================

% Esta sección se apoya en hechos del propio proyecto -plataformas, baterías,
% consumo- y no invoca normativa ambiental concreta, que habría que verificar.

El impacto ambiental del prototipo se concentra en tres frentes: el hardware
empleado, la energía consumida durante la operación y la evaluación, y el
tratamiento de los componentes al final de su vida útil.

\subsection{Reutilización de hardware existente}

Las dos plataformas robóticas empleadas son equipo institucional preexistente,
no una adquisición generada por el proyecto. El presupuesto del anteproyecto las
contabiliza como uso y desgaste, no como compra. Esta decisión, tomada
inicialmente por razones económicas, reduce la huella material del trabajo:
el proyecto no introduce plataformas nuevas al inventario y extiende la vida útil
de equipos que de otro modo permanecerían subutilizados.

Del mismo modo, el entorno de evaluación es el propio edificio de la universidad,
sin construcción de infraestructura dedicada ni modificación de la existente. El
sistema opera sobre la planta tal como está.

\subsection{Consumo energético y autonomía}

Cada vehículo opera con dos fuentes de alimentación independientes, una para
cómputo y otra para tracción. La duración de una corrida de evaluación es del
orden de decenas de segundos, y una sesión completa de campo consume una
fracción de la carga disponible. El consumo agregado del sistema durante toda la
campaña experimental es, por tanto, marginal frente al de la estación de trabajo
empleada para la simulación, que ejecuta dos instancias del simulador de forma
concurrente y sostenida.

Esa asimetría merece señalarse: en un proyecto de robótica de esta escala,
\textbf{el mayor consumo energético no está en los robots sino en la simulación}.
Es también un argumento a favor de la estrategia de verificación adoptada, en la
que la lógica de coordinación se comprueba como función pura, sin simulador, en
milisegundos y sin coste energético apreciable.

\subsection{Ciclo de vida de las baterías}

Las baterías de polímero de litio que alimentan la tracción son el componente de
mayor sensibilidad ambiental del sistema, por su contenido de litio y por su
degradación con los ciclos de carga. Su gestión adecuada al final de la vida útil
corresponde a los canales de disposición de residuos electrónicos de la
institución, y no debe realizarse por vías domésticas.

% VERIFICAR: si la universidad tiene un programa institucional de gestión de
% RAEE (residuos de aparatos eléctricos y electrónicos), conviene nombrarlo aquí
% en concreto. No se afirma que exista sin comprobarlo.

Durante el desarrollo se observó además que el nivel de carga afecta de forma
medible al comportamiento del vehículo: una misma orden de tracción produjo
velocidades distintas en corridas consecutivas, con la batería como primer
sospechoso. Esa dependencia obliga a registrar el nivel de carga antes y después
de cada medición, lo que constituye una buena práctica experimental y, de paso,
desincentiva el uso de baterías degradadas.
